% Capítulo: La misericordia de Dios (generado desde /home/claude/work/14_Misericordia/14_Misericordia_capitulo.tex)
% Texto ancla: Éxodo 34:6–7. Clasificación: Relación con el ser humano; comunicable.

\section{Una palabra para muchas}

En español, ``tiempo'' sirve para preguntar por el clima (``¿qué tiempo hace hoy?'') y por la duración (``¿cuánto tiempo falta?''). Una sola palabra junta conceptos distintos, y el contexto decide cuál. Con la Escritura ocurre lo inverso: varias palabras hebreas y griegas, cada una con su imagen propia, terminaron traducidas con una sola palabra española. Así pasa con ``misericordia''.

El material de estudio de esta serie dice que ``misericordia'' aparece más de 290 veces y traduce seis palabras del original. Un conteo en la RVR1960, sin los títulos de sección que agregan los editores, da 403 apariciones de ``misericordia'' o ``misericordias'' en 385 versículos (331 en el Antiguo Testamento y 72 en el Nuevo); si se suman ``misericordioso'', ``misericordiosa'' y las demás formas, son 432 apariciones en 403 versículos. Detrás de esa única palabra española hay ocho raíces hebreas, dos términos arameos y cuatro familias de palabras griegas. La tabla resume el espectro; la sección~\ref{sec:espectro} lo desarrolla.

\begin{center}\small\renewcommand{\arraystretch}{1.18}
\begin{tabular}{@{}L{5.2cm}L{4.4cm}r@{}}
\toprule
\textbf{Raíz o familia} & \textbf{Sentido base} & \textbf{Versículos} \\ \midrule
\heb{חסד}: \emph{jésed}, \emph{jasid} & amor leal del pacto & $\approx$\,220 \\
\heb{רחם}: \emph{rajamím}, \emph{rajam}, \emph{rajum} & entrañas, compasión & $\approx$\,77 \\
\heb{חנן}: \emph{janán}, \emph{jannun}, \emph{jen}, \emph{tejinná} & gracia; gracia suplicada & $\approx$\,63 \\
\heb{חוס}: \emph{jus} & tener lástima & $\approx$\,11 \\
\heb{חמל}: \emph{jamal}, \emph{jemlá} & perdonar la vida & 10 \\
\heb{נחם}: \emph{nijam} & conmoverse & 1 \\
\heb{טוֹב־עַיִן}: \emph{tov ayin} & ``bueno de ojo'' & 1 \\
Arameo: \emph{rajamín}; \emph{mijan} (de \emph{janán}) & compasión; gracia & 1; 1 \\
\grc{ἔλεος}, \grc{ἐλεέω}, \grc{ἐλεήμων} y sus negativos & compasión activa & 58 \\
\grc{σπλάγχνα}, \grc{σπλαγχνίζομαι} y compuestos & entrañas & 11 \\
\grc{οἰκτιρμός}, \grc{οἰκτίρμων} & compasión & 6 \\
\grc{ὅσιος} (Hch 13:34) & lealtad (= \emph{jésed}) & 1 \\
\bottomrule
\end{tabular}\par
{\footnotesize\color{gris}\textit{Tabla 1.} El espectro de ``misericordia'' en la RVR1960. \textit{Nota.} Conteo propio por versículo con la RVR1960, la BHS y NA27 con números Strong; un versículo puede contener dos raíces.}
\end{center}

\precision{el conteo es propio y aproximado: depende de la versión, de la edición del texto original y de cómo se agrupen las formas; conviene verificarlo en una concordancia (e-Sword). Las ``seis palabras'' del material de estudio corresponden al núcleo: cuatro hebreas (\emph{jésed}, \emph{janán}, \emph{rajam}, \emph{jamal}) y dos griegas (\emph{éleos}, \emph{oiktirmós}). El resto del espectro son términos menos frecuentes, pero cada uno aporta un matiz.}

%===================================================================
\section{Contexto: la gloria cambiada}

\subsection{Babel, la “puerta de dios”}
La revelación de la misericordia en el Sinaí tiene un trasfondo largo: la historia de un hombre que busca su propia gloria. En Babel los hombres dicen: ``edifiquémonos una ciudad y una torre, cuya cúspide llegue al cielo; y hagámonos un nombre'' (Gn 11:4).

\hebverso{עַל־כֵּן קָרָא שְׁמָהּ בָּבֶל כִּי־שָׁם בָּלַל יְהוָה שְׂפַת כָּל־הָאָרֶץ}
\begin{verso}{Génesis 11:9}
Por esto fue llamado el nombre de ella Babel, porque allí confundió \SENOR{} el lenguaje de toda la tierra, y desde allí los esparció sobre la faz de toda la tierra.
\end{verso}

\begin{filologia}{Gn 11:9}
\heb{בָּבֶל} \emph{Bavel} & --- & En acadio, \emph{Bāb-ilim}: \emph{bāb}, ``puerta'', e \emph{ilim}, ``de dios''. Para los babilonios, su ciudad era el punto de contacto entre el cielo y la tierra, el centro del cosmos. \\
\heb{בָּלַל} \emph{balal} & \heb{בלל} & ``confundir, mezclar, revolver''. Moisés hace un juego de palabras: la ``puerta de dios'' es en realidad ``confusión''. \\
\end{filologia}

\precision{la etimología acadia y el juego de palabras con \emph{balal} siguen a Wenham (1987) y a Koehler y Baumgartner (2001). El juego no pretende ser etimología científica: es un comentario teológico del narrador sobre la pretensión de Babilonia.}

\subsection{Cambiaron la gloria}
Pablo describe el mismo movimiento en toda la humanidad:

\begin{verso}{Romanos 1:22–23}
Profesando ser sabios, se hicieron necios, y cambiaron la gloria del Dios incorruptible en semejanza de imagen de hombre corruptible, de aves, de cuadrúpedos y de reptiles.
\end{verso}

Su lenguaje viene del Salmo que recuerda el pecado de Israel en el Sinaí:

\hebverso{יַעֲשׂוּ־עֵגֶל בְּחֹרֵב וַיִּשְׁתַּחֲווּ לְמַסֵּכָה \\ וַיָּמִירוּ אֶת־כְּבוֹדָם בְּתַבְנִית שׁוֹר אֹכֵל עֵשֶׂב}
\begin{verso}{Salmo 106:19–20}
Hicieron becerro en Horeb, Se postraron ante una imagen de fundición. Así cambiaron su gloria Por la imagen de un buey que come hierba.
\end{verso}

\begin{filologia}{Sal 106:20 y Ro 1:23}
\heb{וַיָּמִירוּ} \emph{vayyamiru} & \heb{מור} & ``y cambiaron, trocaron'': cambiar una cosa por otra de menor valor. \\
\heb{כְּבוֹדָם} \emph{kevodam} & \heb{כבד} & ``su gloria'': el Dios que era la gloria de Israel. \emph{Kavod} es literalmente ``peso''. \\
\grc{ἤλλαξαν} \emph{éllaxan} & \grc{ἀλλάσσω} & ``cambiaron'' (Ro 1:23). La LXX del Salmo usa el mismo verbo: \grc{ἠλλάξαντο τὴν δόξαν αὐτῶν} (Sal 105:20 LXX). \\
\end{filologia}

\precision{el eco de Sal 106:20 en Ro 1:23 es reconocido por los comentaristas (Moo, 1996). Para Pablo, lo que hizo Israel ante el becerro es un retrato de lo que hace toda la humanidad.}

\subsection{Abram: elegir la gloria del Dios verdadero}
La historia bíblica traza un arco: de la gloria del Eterno a la gloria de los dioses paganos, y de ahí de vuelta a la gloria del Eterno. El punto de giro es Abram. ``Desde Abram hacia adelante se comienza a llamar a los hombres a elegir la gloria del Dios verdadero.'' Se dan los primeros pasos para formar un pueblo libre de pensamientos errados: un pueblo que tomaría conceptos de los pueblos antiguos, corrigiendo su error y su significado.

\begin{center}
\begin{tikzpicture}[font=\small,node distance=1.0cm]
 \node[draw=dorado,line width=1.6pt,fill=azul,text=papiro,rounded corners=3pt,text width=3.4cm,align=center,inner sep=6pt] (a) {\textbf{Gloria del Eterno}\\ creación\\ \footnotesize Sal 19:1; Ro 1:20};
 \node[draw=azul,line width=0.8pt,fill=gris!12,rounded corners=3pt,text width=3.4cm,align=center,inner sep=6pt,right=of a] (b) {\textbf{Gloria de los dioses paganos}\\ Babel, ídolos, becerro\\ \footnotesize Gn 11; Sal 106:20; Ro 1:23};
 \node[draw=dorado,line width=1.6pt,fill=azul,text=papiro,rounded corners=3pt,text width=3.4cm,align=center,inner sep=6pt,right=of b] (c) {\textbf{Gloria del Eterno}\\ Abram, Sinaí, Cristo\\ \footnotesize Gn 12; Éx 33:18; 2 Co 4:6};
 \draw[-{Latex[length=3mm]},dorado,line width=1.4pt] (a)--(b);
 \draw[-{Latex[length=3mm]},dorado,line width=1.4pt] (b)--(c);
\end{tikzpicture}\par
{\footnotesize\color{gris}\textit{Figura 1.} De la gloria del Eterno a la gloria del Eterno. \textit{Nota.} Elaboración propia a partir del material de estudio de la serie.}
\end{center}

\hebverso{וְאֶעֶשְׂךָ לְגוֹי גָּדוֹל וַאֲבָרֶכְךָ וַאֲגַדְּלָה שְׁמֶךָ וֶהְיֵה בְּרָכָה}
\begin{verso}{Génesis 12:1–2}
Pero \SENOR{} había dicho a Abram: Vete de tu tierra y de tu parentela, y de la casa de tu padre, a la tierra que te mostraré. Y haré de ti una nación grande, y te bendeciré, y engrandeceré tu nombre, y serás bendición.
\end{verso}

En Babel el hombre quería ``hacerse un nombre'' (\heb{שֵׁם}, Gn 11:4); a Abram, Dios le dice: ``engrandeceré tu nombre'' (\heb{וַאֲגַדְּלָה שְׁמֶךָ}). El nombre que el hombre no pudo fabricarse, Dios lo da. Y cuando Abram llega a Canaán, el texto anota dónde estaba:

\begin{verso}{Génesis 12:6–7}
Y pasó Abram por aquella tierra hasta el lugar de Siquem, hasta el encino de More; y el cananeo estaba entonces en la tierra. Y apareció \SENOR{} a Abram, y le dijo: A tu descendencia daré esta tierra. Y edificó allí un altar a \SENOR, quien le había aparecido.
\end{verso}

En medio del culto cananeo, Abram levanta un altar al \SENOR. Es su primer ``no'' a los dioses de la tierra. Siglos después, en ese mismo lugar, Israel dirá ``sí'' a un dios cananeo (sección~\ref{sec:siquem}).

\precision{el contraste ``Babel busca un nombre / Abram lo recibe'' es común entre los comentaristas (Wenham, 1987). La lectura de Gn 12:6--7 como rechazo del culto cananeo de Siquem es propia de esta serie.}

%===================================================================
\section{Texto ancla: Éxodo 34:6--7}

\subsection{Después del becerro}
La revelación más completa del carácter de Dios en el Antiguo Testamento llega en el peor momento de Israel. Recién salidos de Egipto, ya en el pacto, fabrican un becerro de oro y dicen: ``Israel, estos son tus dioses, que te sacaron de la tierra de Egipto'' (Éx 32:4). Moisés rompe las tablas (32:19) e intercede: ``que perdones ahora su pecado, y si no, ráeme ahora de tu libro que has escrito'' (32:32). Y entonces pide algo inaudito:

\hebverso{הַרְאֵנִי נָא אֶת־כְּבֹדֶךָ \\ וְחַנֹּתִי אֶת־אֲשֶׁר אָחֹן וְרִחַמְתִּי אֶת־אֲשֶׁר אֲרַחֵם}
\begin{verso}{Éxodo 33:18–19}
El entonces dijo: Te ruego que me muestres tu gloria. Y le respondió: Yo haré pasar todo mi bien delante de tu rostro, y proclamaré el nombre de \SENOR{} delante de ti; y tendré misericordia del que tendré misericordia, y seré clemente para con el que seré clemente.
\end{verso}

\begin{filologia}{Éx 33:18--19}
\heb{כָּבוֹד} \emph{kavod} & \heb{כבד} & ``gloria'', literalmente ``peso'': Moisés pide ver el peso, la realidad de Dios. \\
\heb{טוּב} \emph{tuv} & \heb{טוב} & ``bien, bondad'': Dios responde al pedido de gloria con ``todo mi bien''. \\
\heb{וְחַנֹּתִי} \emph{vejannotí} & \heb{חנן} & ``tendré gracia'' de quien tendré gracia (RVR1960: ``tendré misericordia''). \\
\heb{וְרִחַמְתִּי} \emph{verijamtí} & \heb{רחם} & ``me compadeceré'' de quien me compadeceré (RVR1960: ``seré clemente''). \\
\end{filologia}

\precision{la RVR1960 invierte los términos: traduce \emph{janán} por ``misericordia'' y \emph{rajam} por ``clemente''. La construcción ``tendré gracia de quien tendré gracia'' (\emph{idem per idem}) es la misma de ``Yo soy el que soy'' (Éx 3:14): afirma la libertad soberana de Dios.}

\subsection{El Nombre proclamado}
Moisés sube de nuevo al Sinaí con dos tablas nuevas (Éx 34:4). Dios desciende en la nube y pasa delante de él:

\hebverso{וַיַּעֲבֹר יְהוָה עַל־פָּנָיו וַיִּקְרָא יְהוָה יְהוָה אֵל רַחוּם וְחַנּוּן \\ אֶרֶךְ אַפַּיִם וְרַב־חֶסֶד וֶאֱמֶת נֹצֵר חֶסֶד לָאֲלָפִים \\ נֹשֵׂא עָוֹן וָפֶשַׁע וְחַטָּאָה וְנַקֵּה לֹא יְנַקֶּה}
\begin{verso}{Éxodo 34:6–7}
Y pasando \SENOR{} por delante de él, proclamó: ¡\SENOR! ¡\SENOR! fuerte, misericordioso y piadoso; tardo para la ira, y grande en misericordia y verdad; que guarda misericordia a millares, que perdona la iniquidad, la rebelión y el pecado, y que de ningún modo tendrá por inocente al malvado; que visita la iniquidad de los padres sobre los hijos y sobre los hijos de los hijos, hasta la tercera y cuarta generación.
\end{verso}

\begin{filologia}{Éx 34:6--7}
\heb{אֵל} \emph{El} & --- & ``Dios'' (RVR1960: ``fuerte''). \\
\heb{רַחוּם} \emph{rajum} & \heb{רחם} & ``compasivo'', de \emph{réjem}, el vientre materno: la ternura de una madre por su hijo. \\
\heb{חַנּוּן} \emph{jannun} & \heb{חנן} & ``clemente, lleno de gracia'': el que se inclina con favor hacia quien no lo merece. \\
\heb{אֶרֶךְ אַפַּיִם} \emph{érej appáyim} & \heb{ארך}, \heb{אנף} & ``largo de narices'': lento para que la ira se encienda. \\
\heb{רַב־חֶסֶד וֶאֱמֶת} \emph{rav jésed ve'emet} & \heb{חסד}, \heb{אמן} & ``abundante en amor leal y en verdad'': los dos pilares del pacto. \\
\heb{נֹצֵר חֶסֶד לָאֲלָפִים} \emph{notser jésed la'alafim} & \heb{נצר} & ``que guarda (vigila, custodia) el amor leal para millares''. \\
\heb{נֹשֵׂא עָוֹן} \emph{nosé avón} & \heb{נשׂא} & ``que levanta, carga la iniquidad'': perdonar es cargar con la culpa del otro. \\
\heb{וְנַקֵּה לֹא יְנַקֶּה} \emph{venaqqé lo yenaqqé} & \heb{נקה} & ``y absolver, no absuelve'': el infinitivo absoluto refuerza la negación (RVR1960: ``de ningún modo tendrá por inocente''). \\
\end{filologia}

El Nombre tiene tres partes. Primero, \emph{lo que Dios es}: compasivo, clemente, lento para la ira, abundante en \emph{jésed} y verdad (v.~6). Segundo, \emph{lo que hace}: guarda el \emph{jésed} a millares y carga la iniquidad, la rebelión y el pecado (v.~7a). Tercero, \emph{lo que no hace}: no absuelve al culpable impenitente (v.~7b). Las proporciones hablan por sí mismas: \emph{jésed} para ``millares'' (de generaciones, como entiende Dt 7:9) frente a visitación hasta la tercera y cuarta. La misericordia es más grande que el juicio, sin anularlo.

\begin{verso}{Éxodo 34:8}
Entonces Moisés, apresurándose, bajó la cabeza hacia el suelo y adoró.
\end{verso}

Moisés había pedido ver la \emph{gloria} (33:18), y Dios le mostró su \emph{carácter}. Esa es la tesis de esta obra: la gloria de Dios es la comunión de sus atributos. En el Sinaí, esa comunión se anuncia con la misericordia en primer lugar.

\subsection{La fórmula que recorre la Biblia}
Israel no inventó un concepto nuevo de Dios en cada época: repitió el Nombre que Dios proclamó en el Sinaí. La fórmula de Éx 34:6 reaparece en la oración, en la alabanza y en la profecía.

\begin{center}\small\renewcommand{\arraystretch}{1.2}
\begin{tabular}{@{}L{2.3cm}L{8.8cm}L{3.6cm}@{}}
\toprule
\textbf{Texto} & \textbf{RVR1960} & \textbf{Uso} \\ \midrule
Nm 14:18--19 & ``\SENOR, tardo para la ira y grande en misericordia\ldots{} Perdona ahora la iniquidad de este pueblo según la grandeza de tu misericordia'' & Moisés intercede \\
Sal 86:15 & ``Mas tú, Señor, Dios misericordioso y clemente, lento para la ira, y grande en misericordia y verdad'' & oración personal \\
Sal 103:8 & ``Misericordioso y clemente es \SENOR; lento para la ira, y grande en misericordia'' & alabanza \\
Sal 145:8 & ``Clemente y misericordioso es \SENOR, lento para la ira, y grande en misericordia'' & alabanza (acróstico) \\
Jl 2:13 & ``porque misericordioso es y clemente, tardo para la ira y grande en misericordia'' & llamado al arrepentimiento \\
Jon 4:2 & ``sabía yo que tú eres Dios clemente y piadoso, tardo en enojarte, y de grande misericordia'' & queja del profeta \\
Nah 1:3 & ``\SENOR{} es tardo para la ira\ldots{} y no tendrá por inocente al culpable'' & juicio de Nínive \\
\bottomrule
\end{tabular}\par
{\footnotesize\color{gris}\textit{Tabla 2.} Ecos de Éx 34:6--7 en el Antiguo Testamento.}
\end{center}

%===================================================================
\section{La palabra: cuatro columnas}

La Escritura dice la misericordia con muchas palabras, pero cuatro de ellas, las del material de estudio, sostienen el techo. Ninguna basta sola; juntas muestran cómo Dios se acerca al que no merece (\emph{janán}), al que pertenece a su pacto (\emph{jésed}), al que sufre (\emph{rajamím}) y al que debería morir (\emph{jamal}).

\subsection{El espectro completo}\label{sec:espectro}
Antes de las cuatro columnas conviene ver el conjunto. Estos son todos los términos del texto original que la RVR1960 traduce con ``misericordia'', ``misericordias'', ``misericordioso'' o ``tener misericordia''.

\begin{center}\small\renewcommand{\arraystretch}{1.2}
\begin{tabular}{@{}L{1.3cm}L{4.6cm}L{4.4cm}L{4.0cm}@{}}
\toprule
\textbf{Raíz} & \textbf{Términos} & \textbf{Imagen y sentido} & \textbf{Textos} \\ \midrule
\heb{חסד} & \heb{חֶסֶד} \emph{jésed}; \heb{חָסִיד} \emph{jasid} ``misericordioso''; \heb{הִתְחַסֵּד} \emph{hitjassed} ``mostrarse misericordioso'' & lealtad del pacto; el \emph{jasid} es el que vive en \emph{jésed} & Éx 34:6; Sal 136; 2~S 22:26; Jer 3:12; Mi 7:2 \\
\heb{רחם} & \heb{רַחֲמִים} \emph{rajamím}; \heb{רָחַם} \emph{rajam}; \heb{רַחוּם} \emph{rajum} & el vientre (\emph{réjem}): ternura de entrañas & Sal 103:13; Jer 31:20; Os 2:23 \\
\heb{חנן} & \heb{חָנַן} \emph{janán}; \heb{חַנּוּן} \emph{jannun}; \heb{חֵן} \emph{jen}; \heb{תְּחִנָּה} \emph{tejinná}; \heb{תַּחֲנוּנִים} \emph{tajanunim} & inclinarse con favor; la gracia que se suplica & Sal 86:3; Pr 14:21; Jos 11:20; Esd 9:8; Jer 31:9 \\
\heb{חמל} & \heb{חָמַל} \emph{jamal}; \heb{חֶמְלָה} \emph{jemlá} & perdonar la vida, no destruir & Gn 19:16; Mal 3:17; Ez 5:11 \\
\heb{חוס} & \heb{חוּס} \emph{jus} & el ojo que ``tiene lástima'' & Sal 72:13; Jon 4:11; Ez 5:11 \\
\heb{נחם} & \heb{נִחַם} \emph{nijam} (nifal) & conmoverse hasta cambiar el curso & Jue 2:18 \\
\heb{טוב} & \heb{טוֹב־עַיִן} \emph{tov ayin} & ``bueno de ojo'': el generoso & Pr 22:9 \\
Arameo & \heb{רַחֲמִין} \emph{rajamín}; \heb{מִחַן} \emph{mijan} (de \heb{חֲנַן}) & compasión; hacer gracia al oprimido & Dn 2:18; 4:27 \\
\bottomrule
\end{tabular}\par
{\footnotesize\color{gris}\textit{Tabla 3.} Términos hebreos y arameos traducidos ``misericordia'' en la RVR1960. \textit{Nota.} Sentidos según Strong y Koehler y Baumgartner (2001).}
\end{center}

\begin{center}\small\renewcommand{\arraystretch}{1.2}
\begin{tabular}{@{}L{3.6cm}L{4.8cm}L{1.3cm}L{4.6cm}@{}}
\toprule
\textbf{Término} & \textbf{Sentido} & \textbf{Vv.} & \textbf{Textos} \\ \midrule
\grc{ἔλεος}, \grc{ἐλεέω} & misericordia; tener misericordia (en la LXX traduce \emph{jésed}) & 26; 28 & Mt 9:13; Ro 9:15; Ef 2:4 \\
\grc{ἐλεήμων} & misericordioso & 2 & Mt 5:7; He 2:17 \\
\grc{ἀνέλεος}, \grc{ἀνελεήμων} & sin misericordia & 1; 1 & Stg 2:13; Ro 1:31 \\
\grc{οἰκτιρμός}, \grc{οἰκτίρμων} & compasión; compasivo (LXX: \emph{rajamím}, \emph{rajum}) & 4; 2 & Ro 12:1; 2~Co 1:3; Lc 6:36 \\
\grc{σπλαγχνίζομαι} & conmoverse en las entrañas & 5 & Mt 18:27; Lc 10:33; 15:20 \\
\grc{σπλάγχνα} & ``entrañas'': ``entrañable misericordia'' & 3 & Lc 1:78; Flp 2:1; Col 3:12 \\
\grc{εὔσπλαγχνος} & ``de buenas entrañas'' & 2 & Ef 4:32; 1~P 3:8 \\
\grc{πολύσπλαγχνος} & ``de muchas entrañas'' & 1 & Stg 5:11 \\
\grc{ὅσιος} & \grc{τὰ ὅσια Δαυίδ}: las lealtades prometidas a David & 1 & Hch 13:34 (= Is 55:3, \emph{jasdé David}) \\
\bottomrule
\end{tabular}\par
{\footnotesize\color{gris}\textit{Tabla 4.} Términos griegos traducidos ``misericordia'' en la RVR1960. \textit{Nota.} Versículos según NA27 con números Strong; sentidos según BDAG.}
\end{center}

\precision{dos casos especiales del Nuevo Testamento: en Mt 21:41 ``sin misericordia'' traduce el adverbio \grc{κακῶς} (``malamente''), una traducción libre; en Tit 1:4, ``misericordia'' (\grc{ἔλεος}) está en el Texto Recibido que sigue la RVR1960, pero no en NA28. En hebreo, \emph{nijam} suele traducirse ``arrepentirse'' (de Dios: cambiar el curso anunciado); solo en Jue 2:18 la RVR1960 lo vierte ``era movido a misericordia''.}

\subsubsection*{Las palabras menores}
Cada término menor añade un matiz. El \emph{jasid} es el hombre en quien el \emph{jésed} se volvió carácter: ``Faltó el misericordioso de la tierra, y ninguno hay recto entre los hombres'' (Mi 7:2). Dios lo dice de sí mismo: ``misericordioso (\heb{חָסִיד}) soy yo, dice \SENOR, no guardaré para siempre el enojo'' (Jer 3:12). Y el salmo enseña que Dios responde en la misma moneda: ``Con el misericordioso te mostrarás misericordioso'' (\heb{עִם־חָסִיד תִּתְחַסָּד}, 2~S 22:26; Sal 18:25).

\emph{Tejinná} y \emph{tajanunim}, de la raíz \emph{janán}, son la gracia vista desde quien la pide. Esdras confiesa: ``por un breve momento ha habido misericordia (\heb{תְחִנָּה}) de parte de \SENOR{} nuestro Dios, para hacer que nos quedase un remanente libre'' (Esd 9:8). Jeremías dice del regreso de Efraín: ``Irán con lloro, mas con misericordia (\heb{וּבְתַחֲנוּנִים}, ``con súplicas'') los haré volver'' (Jer 31:9).

\emph{Jus} es el ojo que se apiada: ``Tendrá misericordia (\heb{יָחֹס}) del pobre y del menesteroso'' (Sal 72:13); en el juicio se dice lo contrario: ``mi ojo no perdonará (\heb{תָחוֹס}), ni tampoco tendré yo misericordia (\heb{אֶחְמוֹל})'' (Ez 5:11), donde \emph{jus} y \emph{jamal} van juntos. \emph{Nijam} describe a Dios conmovido por el dolor de su pueblo: ``\SENOR{} era movido a misericordia (\heb{יִנָּחֵם}) por sus gemidos'' (Jue 2:18). Y \emph{tov ayin}, ``bueno de ojo'', llama misericordioso al generoso: ``El ojo misericordioso será bendito, porque dio de su pan al indigente'' (Pr 22:9).

En arameo, Daniel y sus compañeros piden ``misericordias'' (\heb{רַחֲמִין}) al Dios del cielo (Dn 2:18), y Daniel aconseja a Nabucodonosor redimir sus iniquidades ``haciendo misericordias para con los oprimidos'' (\heb{בְּמִחַן עֲנָיִן}, Dn 4:27).

En griego, la familia de \grc{σπλάγχνα} (``entrañas'') traduce la imagen de \emph{rajamím}: Zacarías canta ``la entrañable misericordia de nuestro Dios'' (\grc{διὰ σπλάγχνα ἐλέους θεοῦ ἡμῶν}, Lc 1:78); Pablo pide vestirse ``de entrañable misericordia'' (\grc{σπλάγχνα οἰκτιρμοῦ}, Col 3:12) y ser ``misericordiosos'' (\grc{εὔσπλαγχνοι}, ``de buenas entrañas'', Ef 4:32); Santiago confiesa que ``el Señor es muy misericordioso (\grc{πολύσπλαγχνος}) y compasivo'' (Stg 5:11), eco de Éx 34:6. Hebreos llama a Cristo ``misericordioso y fiel sumo sacerdote'' (\grc{ἐλεήμων καὶ πιστός}, He 2:17): el par \emph{jésed ve'emet} hecho persona. Y en Antioquía de Pisidia, Pablo anuncia la resurrección con Is 55:3: ``Os daré las misericordias fieles de David'' (\grc{τὰ ὅσια Δαυὶδ τὰ πιστά}, Hch 13:34), donde la LXX traduce \heb{חַסְדֵי דָוִד הַנֶּאֱמָנִים}. El \emph{jésed} prometido a David se cumple en el Resucitado.

\precision{colosenses y Efesios son de autoría paulina discutida; Hebreos no es de Pablo; Hch 13:34 recoge un discurso de Pablo según Lucas.}

De todo el espectro, cuatro palabras son las columnas que sostienen el techo.

\subsection{\emph{Janán}: gracia, favor inmerecido}
\hebverso{יָאֵר יְהוָה פָּנָיו אֵלֶיךָ וִיחֻנֶּךָּ}
\begin{verso}{Números 6:25}
\SENOR{} haga resplandecer su rostro sobre ti, y tenga de ti misericordia.
\end{verso}

\hebverso{חָנֵּנִי אֲדֹנָי כִּי אֵלֶיךָ אֶקְרָא כָּל־הַיּוֹם}
\begin{verso}{Salmo 86:3}
Ten misericordia de mí, oh Señor; Porque a ti clamo todo el día.
\end{verso}

\begin{filologia}{Nm 6:25 y Sal 86:3}
\heb{וִיחֻנֶּךָּ} \emph{viyjunnekka} & \heb{חנן} & ``y te otorgue gracia''. \emph{Janán} (Strong H2603) es ``inclinarse en bondad hacia un inferior''. \\
\heb{חָנֵּנִי} \emph{jonnení} & \heb{חנן} & ``ten gracia de mí'': imperativo del que no tiene derechos que reclamar. \\
\heb{חֵן}, \heb{תְּחִנָּה} \emph{jen}, \emph{tejinná} & \heb{חנן} & ``gracia'' y ``súplica''. La forma reflexiva del verbo (\emph{hitjanén}) significa \emph{implorar}: pedir gracia. \\
\end{filologia}

En la bendición sacerdotal, la gracia va unida al rostro de Dios que resplandece. En el Salmo, el que no tiene méritos clama ``todo el día''. Y en Éx 33:19 Dios afirma que su gracia es libre: ``tendré gracia de quien tendré gracia''.

La libertad de esa gracia tiene un eco en el Nuevo Testamento. Cuando Pedro, después de ser restaurado, pregunta por el destino de Juan, Jesús le responde:

\begin{verso}{Juan 21:21–22}
Cuando Pedro le vio, dijo a Jesús: Señor, ¿y qué de éste? Jesús le dijo: Si quiero que él quede hasta que yo venga, ¿qué a ti? Sígueme tú.
\end{verso}

\precision{la conexión entre Éx 33:19 y Jn 21:22 es del autor de la serie y es una lectura propia: la gracia de Dios no se administra con criterio humano; a cada uno le toca recibirla y seguir.}

La parábola del siervo que no perdonó muestra la gracia en su forma más pura:

\begin{verso}{Mateo 18:26–27}
Entonces aquel siervo, postrado, le suplicaba, diciendo: Señor, ten paciencia conmigo, y yo te lo pagaré todo. El señor de aquel siervo, movido a misericordia, le soltó y le perdonó la deuda.
\end{verso}

Diez mil talentos (Mt 18:24) era una deuda imposible de pagar. El siervo implora, y el rey no le da más plazo: le perdona todo. El verbo es \grc{σπλαγχνισθείς}, ``conmovido en las entrañas''.

\subsection{\emph{Rajamím}: compasión desde las entrañas}
\hebverso{כְּרַחֵם אָב עַל־בָּנִים רִחַם יְהוָה עַל־יְרֵאָיו}
\begin{verso}{Salmo 103:13}
Como el padre se compadece de los hijos, Se compadece \SENOR{} de los que le temen.
\end{verso}

\begin{filologia}{Sal 103:13}
\heb{רֶחֶם} \emph{réjem} & \heb{רחם} & ``vientre, matriz'' (H7358). \\
\heb{רַחֲמִים} \emph{rajamím} & \heb{רחם} & plural de \emph{réjem}: ``entrañas'' y, por extensión, ``compasión'' (H7356). Es el amor que se siente en el cuerpo, como el de la madre de 1~R 3:26, a quien ``sus entrañas se le conmovieron por su hijo''. \\
\heb{כְּרַחֵם}\ldots{} \heb{רִחַם} \emph{kerajem\ldots{} rijam} & \heb{רחם} & ``como se compadece\ldots{} se compadece'': el paralelismo compara la compasión del padre con la de Dios. \\
\end{filologia}

El Salmo aplica a un padre una imagen que nace del vientre materno: la ternura de Dios reúne la de la madre y la del padre. Jeremías la lleva más lejos:

\hebverso{הֲבֵן יַקִּיר לִי אֶפְרַיִם אִם יֶלֶד שַׁעֲשֻׁעִים \\ עַל־כֵּן הָמוּ מֵעַי לוֹ רַחֵם אֲרַחֲמֶנּוּ נְאֻם־יְהוָה}
\begin{verso}{Jeremías 31:18–20}
Escuchando, he oído a Efraín que se lamentaba: Me azotaste, y fui castigado como novillo indómito; conviérteme, y seré convertido, porque tú eres \SENOR{} mi Dios. Porque después que me aparté tuve arrepentimiento, y después que reconocí mi falta, herí mi muslo; me avergoncé y me confundí, porque llevé la afrenta de mi juventud. ¿No es Efraín hijo precioso para mí? ¿no es niño en quien me deleito? pues desde que hablé de él, me he acordado de él constantemente. Por eso mis entrañas se conmovieron por él; ciertamente tendré de él misericordia, dice \SENOR.
\end{verso}

\begin{filologia}{Jer 31:20}
\heb{הָמוּ מֵעַי} \emph{hamú me'ay} & \heb{המה} & ``mis entrañas se agitaron'': el verbo describe el bramido del mar y el ruido de una multitud. \\
\heb{רַחֵם אֲרַחֲמֶנּוּ} \emph{rajem arajamennu} & \heb{רחם} & ``compadecer, me compadeceré de él'': el infinitivo absoluto intensifica (RVR1960: ``ciertamente''). \\
\end{filologia}

Efraín, el reino del norte, había sido castigado ``como novillo indómito''; la palabra es \heb{עֵגֶל} (\emph{éguel}), la misma del becerro de Éx 32. El hijo vuelve avergonzado, y el Padre confiesa que sus entrañas se agitan por él. Jesús contó la misma escena:

\begin{verso}{Lucas 15:20}
Y levantándose, vino a su padre. Y cuando aún estaba lejos, lo vio su padre, y fue movido a misericordia, y corrió, y se echó sobre su cuello, y le besó.
\end{verso}

El verbo \grc{ἐσπλαγχνίσθη} (\emph{esplanjnísthē}) viene de \grc{σπλάγχνα}, ``entrañas'': es el equivalente griego de \emph{rajamím}. El padre de la parábola actúa como el Dios de Jer 31:20. Y Jesús mismo se conmueve ante el dolor: ante la tumba de Lázaro ``se estremeció en espíritu y se conmovió\ldots{} Jesús lloró'' (Jn 11:33, 35).

\precision{jn 11:33 no usa \grc{σπλαγχνίζομαι}, sino \grc{ἐνεβριμήσατο} (``se estremeció'') y \grc{ἐτάραξεν ἑαυτόν} (``se turbó''); el v.~35 dice \grc{ἐδάκρυσεν} (``lloró''). El verbo de las entrañas aplicado a Jesús aparece en los sinópticos (Mt 9:36; 14:14; Lc 7:13).}

\subsection{\emph{Jamal}: perdonar la vida, levantar}
\hebverso{בְּחֶמְלַת יְהוָה עָלָיו}
\begin{verso}{Génesis 19:16}
Y deteniéndose él, los varones asieron de su mano, y de la mano de su mujer y de las manos de sus dos hijas, según la misericordia de \SENOR{} para con él; y lo sacaron y lo pusieron fuera de la ciudad.
\end{verso}

\hebverso{וְחָמַלְתִּי עֲלֵיהֶם כַּאֲשֶׁר יַחְמֹל אִישׁ עַל־בְּנוֹ הָעֹבֵד אֹתוֹ}
\begin{verso}{Malaquías 3:17}
Y serán para mí especial tesoro, ha dicho \SENOR{} de los ejércitos, en el día en que yo actúe; y los perdonaré, como el hombre que perdona a su hijo que le sirve.
\end{verso}

\begin{filologia}{Gn 19:16 y Mal 3:17}
\heb{חֶמְלָה} \emph{jemlá} & \heb{חמל} & ``compasión, piedad'' (Gn 19:16): la que libra de la destrucción. \\
\heb{וְחָמַלְתִּי} \emph{vejamaltí} & \heb{חמל} & ``y tendré piedad, los perdonaré'' (Mal 3:17). \emph{Jamal} (H2550) significa compadecerse y, por implicación, librar, perdonar la vida, no destruir. \\
\end{filologia}

\emph{Jamal} es la misericordia que \emph{saca} del lugar del juicio. Lot ``se demoraba''; la misericordia lo toma de la mano y lo pone fuera de Sodoma. En Malaquías, Dios ``perdona'' a los suyos en el día del juicio como un padre que preserva a su hijo.

Jonás termina con una pregunta de Dios que usa un sinónimo cercano:

\hebverso{וַאֲנִי לֹא אָחוּס עַל־נִינְוֵה הָעִיר הַגְּדוֹלָה}
\begin{verso}{Jonás 4:10–11}
Y dijo \SENOR: Tuviste tú lástima de la calabacera, en la cual no trabajaste, ni tú la hiciste crecer; que en espacio de una noche nació, y en espacio de otra noche pereció. ¿Y no tendré yo piedad de Nínive, aquella gran ciudad donde hay más de ciento veinte mil personas que no saben discernir entre su mano derecha y su mano izquierda, y muchos animales?
\end{verso}

\precision{jon 4:10--11 no usa \emph{jamal}, sino \heb{חוּס} (\emph{jus}, H2347), ``tener lástima, perdonar la vida''. Jonás ``tuvo lástima'' de una planta; Dios, de una ciudad entera y hasta de sus animales. El libro termina con la pregunta sin respuesta: la deja al lector.}

\subsection{\emph{Jésed}: amor leal del pacto}
\hebverso{חַסְדֵי יְהוָה כִּי לֹא־תָמְנוּ כִּי לֹא־כָלוּ רַחֲמָיו \\ חֲדָשִׁים לַבְּקָרִים רַבָּה אֱמוּנָתֶךָ}
\begin{verso}{Lamentaciones 3:22–23}
Por la misericordia de \SENOR{} no hemos sido consumidos, porque nunca decayeron sus misericordias. Nuevas son cada mañana; grande es tu fidelidad.
\end{verso}

\hebverso{כָּל־אָרְחוֹת יְהוָה חֶסֶד וֶאֱמֶת לְנֹצְרֵי בְרִיתוֹ וְעֵדֹתָיו}
\begin{verso}{Salmo 25:10}
Todas las sendas de \SENOR{} son misericordia y verdad, Para los que guardan su pacto y sus testimonios.
\end{verso}

\begin{filologia}{Lam 3:22 y Sal 25:10}
\heb{חֶסֶד} \emph{jésed} & \heb{חסד} & bondad leal, amor que se mantiene por compromiso (H2617). Es la palabra del pacto: la lealtad que Dios jura y cumple. \\
\heb{חַסְדֵי} \emph{jasdé} & \heb{חסד} & ``las misericordias (de)'': plural de \emph{jésed}; junto a \heb{רַחֲמָיו} \emph{rajamav}, ``sus compasiones'' (Lam 3:22). \\
\heb{בְּרִית} \emph{berit} & --- & ``pacto'' (Sal 25:10): el \emph{jésed} se vive dentro de un pacto. \\
\end{filologia}

Lamentaciones 3 es un acróstico triple (tres versos por cada letra del alfabeto), y estos versos están en el centro del libro, escritos entre las ruinas de Jerusalén. En medio del juicio, el poeta encuentra que el \emph{jésed} no se agotó.

\precision{\emph{Jésed} no tiene equivalente exacto en español. La RVR1960 lo traduce ``misericordia'' 205 veces, pero también ``bondad'', ``benignidad'', ``favor'' y ``fidelidad'' (Strong H2617).}

En Oseas, \emph{jésed} es lo que Dios espera de su pueblo:

\hebverso{כִּי חֶסֶד חָפַצְתִּי וְלֹא־זָבַח וְדַעַת אֱלֹהִים מֵעֹלוֹת}
\begin{verso}{Oseas 6:6}
Porque misericordia quiero, y no sacrificio, y conocimiento de Dios más que holocaustos.
\end{verso}

Dos versos antes, Dios se queja: ``La piedad vuestra (\emph{jésed}) es como nube de la mañana'' (Os 6:4). Jesús cita Os 6:6 dos veces (Mt 9:13; 12:7). La LXX traduce \emph{jésed} con \grc{ἔλεος}: \grc{διότι ἔλεος θέλω καὶ οὐ θυσίαν}. El puente entre los testamentos es directo: \emph{jésed} $\rightarrow$ \emph{éleos}.

\subsection{Del hebreo a Pablo}
La Septuaginta tradujo el Nombre de Éx 34:6 así: \grc{κύριος ὁ θεὸς οἰκτίρμων καὶ ἐλεήμων, μακρόθυμος καὶ πολυέλεος καὶ ἀληθινός}. Ese vocabulario pasa a Pablo.

\begin{center}
\begin{tikzpicture}[font=\small,node distance=0.8cm]
 \node[draw=azul,line width=0.8pt,fill=papiro,rounded corners=3pt,text width=3.6cm,align=center,inner sep=5pt] (h) {\textbf{Hebreo (Éx 34:6)}\\ \heb{רַחוּם} \emph{rajum}\\ \heb{חַנּוּן} \emph{jannun}\\ \heb{רַב־חֶסֶד} \emph{rav jésed}};
 \node[draw=azul,line width=0.8pt,fill=papiro,rounded corners=3pt,text width=3.6cm,align=center,inner sep=5pt,right=of h] (l) {\textbf{Septuaginta}\\ \grc{οἰκτίρμων}\\ \grc{ἐλεήμων}\\ \grc{πολυέλεος}};
 \node[draw=azul,line width=0.8pt,fill=papiro,rounded corners=3pt,text width=4.2cm,align=center,inner sep=5pt,right=of l] (p) {\textbf{Pablo}\\ \grc{πατὴρ τῶν οἰκτιρμῶν} (2 Co 1:3)\\ \grc{πλούσιος ἐν ἐλέει} (Ef 2:4)\\ \grc{ἐλεήσω ὃν ἂν ἐλεῶ} (Ro 9:15)};
 \draw[-{Latex[length=2.5mm]},dorado,line width=1.2pt] (h)--(l);
 \draw[-{Latex[length=2.5mm]},dorado,line width=1.2pt] (l)--(p);
\end{tikzpicture}\par
{\footnotesize\color{gris}\textit{Figura 2.} Del Nombre de Éx 34:6 al vocabulario de Pablo. \textit{Nota.} Elaboración propia (BHS, LXX, NA28).}
\end{center}

\begin{filologia}{El vocabulario griego}
\grc{ἔλεος} \emph{éleos} & --- & compasión activa ante la miseria; en la LXX traduce casi siempre \emph{jésed} (BDAG). \\
\grc{οἰκτιρμός} \emph{oiktirmós} & \grc{οἰκτίρω} & compasión, piedad que se conmueve; en la LXX traduce sobre todo \emph{rajamím}. \\
\grc{σπλαγχνίζομαι} \emph{splanjnízomai} & \grc{σπλάγχνα} & conmoverse en las entrañas; solo en los evangelios. \\
\end{filologia}

%===================================================================
\section{Definición}

\begin{tcolorbox}[colback=papiro,colframe=azul,boxrule=1pt,arc=3pt,fonttitle=\bfseries,title=Definición]
La misericordia de Dios es la perfección eterna por la cual Dios, movido desde sus entrañas y fiel a su pacto, se inclina hacia el miserable y el culpable para aliviarlo, perdonarlo y levantarlo, sin dejar de ser justo (Éx 34:6--7).
\end{tcolorbox}

\begin{center}\small\renewcommand{\arraystretch}{1.2}
\begin{tabular}{@{}L{3.4cm}L{6.6cm}L{4.4cm}@{}}
\toprule
\textbf{Rasgo} & \textbf{Qué afirma} & \textbf{Textos} \\ \midrule
Libre (\emph{janán}) & no se merece ni se exige & Éx 33:19; Ro 9:15 \\
Leal (\emph{jésed}) & se apoya en el pacto & Sal 25:10; Lam 3:22 \\
Tierna (\emph{rajamím}) & nace de las entrañas & Sal 103:13; Jer 31:20 \\
Que salva (\emph{jamal}) & saca del lugar del juicio & Gn 19:16; Mal 3:17 \\
Justa & no absuelve al impenitente & Éx 34:7; Nm 14:18 \\
\bottomrule
\end{tabular}\end{center}

\textbf{Clasificación.} La misericordia pertenece a los atributos de la \emph{relación con el ser humano}: responde a la pregunta ``¿cómo se revela Dios?''. Es eterna como todas sus perfecciones (``para siempre es su misericordia'', Sal 136:1), aunque se ejerce sobre criaturas que sufren. Es un atributo \emph{comunicable}: ``Sed, pues, misericordiosos, como también vuestro Padre es misericordioso'' (Lc 6:36).

\subsection{Misericordia y gracia}
Pink explica que la misericordia tiene su origen en la bondad de Dios. La bondad da a la criatura por el solo hecho de ser criatura; la misericordia es ``la pronta inclinación de Dios a aliviar la miseria de las criaturas caídas. Así, la misericordia presupone el pecado'' (Pink, 2020, cap. 14). Para distinguirla de la gracia, Pink recurre a los ángeles que no cayeron: Dios los eligió y los preservó, de modo que reciben gracia; pero nunca han recibido misericordia, porque nunca la necesitaron.

La distinción puede resumirse así: la gracia mira al \emph{culpable} que no merece; la misericordia mira al \emph{miserable} que sufre. En el hombre caído ambas cosas coinciden, y por eso Hebreos las junta: ``Acerquémonos, pues, confiadamente al trono de la gracia, para alcanzar misericordia y hallar gracia para el oportuno socorro'' (He 4:16).

\subsection{Siempre juntas: \emph{jésed} y \emph{emet}}
\hebverso{חֶסֶד־וֶאֱמֶת נִפְגָּשׁוּ צֶדֶק וְשָׁלוֹם נָשָׁקוּ}
\begin{verso}{Salmo 85:10 [85:11]}
La misericordia y la verdad se encontraron; La justicia y la paz se besaron.
\end{verso}

\begin{filologia}{Sal 85:10 (85:11)}
\heb{אֱמֶת} \emph{emet} & \heb{אמן} & ``firmeza, fidelidad, verdad'': lo que hace confiable el \emph{jésed}. \\
\heb{נִפְגָּשׁוּ} \emph{nifgashu} & \heb{פגשׁ} & ``se encontraron'', como dos personas que salen una al encuentro de la otra. \\
\end{filologia}

El par \emph{jésed ve'emet} aparece desde Éx 34:6 y recorre la Escritura: ``todas las sendas de \SENOR{} son misericordia y verdad'' (Sal 25:10); ``cumplirás la verdad a Jacob, y a Abraham la misericordia'' (Mi 7:20). Juan describe al Verbo encarnado ``lleno de gracia y de verdad'' (Jn 1:14). Misericordia sin verdad sería sentimentalismo; verdad sin misericordia, dureza. En Dios van siempre juntas.

\precision{leer ``lleno de gracia y de verdad'' (Jn 1:14) como eco de \emph{rav jésed ve'emet} es una lectura extendida entre los comentaristas (Carson, 1991).}

%===================================================================
\section{Teología: la exposición de Pink}

Pink dedica dos capítulos a este atributo: ``La misericordia de Dios'' y ``La tierna misericordia de Dios'' (Pink, 2020, caps. 14--15). Se usa aquí como referencia, sin seguir su orden.

\begin{verso}{Salmo 136:1}
Alabad a \SENOR, porque él es bueno; porque para siempre es su misericordia.
\end{verso}

Pink abre con el Salmo 136, que repite ``porque para siempre es su misericordia'' en sus 26 versos: la creación, el éxodo, el desierto y la entrada en la tierra se leen como obras del \emph{jésed}. Sus ideas principales son cinco:

\begin{enumerate}
\item La misericordia se origina en la bondad divina y presupone la miseria del pecado.
\item Hay que distinguir tres niveles: una misericordia \emph{general}, sobre toda la creación (``sus misericordias sobre todas sus obras'', Sal 145:9); una misericordia \emph{especial} hacia todos los hombres, a pesar de sus pecados (Mt 5:45); y una misericordia \emph{soberana}, reservada para los herederos de la salvación y comunicada en forma de pacto por medio del Mediador.
\item Dios nunca deja de ser misericordioso, porque es una cualidad de su esencia (Sal 116:5), pero el \emph{ejercicio} de su misericordia lo regula su voluntad soberana (Ro 9:15). Hablar de ``merecer'' misericordia es una contradicción (Tit 3:5).
\item Es una presunción ``diabólica'' seguir pecando contando con la misericordia, porque Dios ha dicho que ``de ningún modo dará por inocente al malvado'' (Éx 34:7). ``Dios muestra misericordia a los verdaderamente penitentes, pero no a los impenitentes'' (Pink, 2020, cap. 14).
\item La ``tierna misericordia'' (\emph{jésed}) de Dios hacia su pueblo está unida a su verdad y a su pacto: ``Mas no quitaré de él mi misericordia, ni falsearé mi verdad. No olvidaré mi pacto'' (Sal 89:33--34). Por eso es motivo de adoración, consuelo en el abatimiento y argumento en la oración (Sal 51:1; 119:76).
\end{enumerate}

\precision{resumen propio de Pink (2020, caps. 14--15). Pink observa también que ningún pueblo antiguo atribuyó a sus dioses una ternura semejante; esa observación coincide con el arco de la sección 2 (de la gloria de los dioses a la gloria del Eterno), pero es un juicio general suyo.}

%===================================================================
\section{Un encuentro en la historia de Israel: Moisés en el Sinaí}

El encuentro que define este atributo es el del texto ancla, leído ahora como historia. Moisés intercede tres veces. La primera, recién bajado del monte, se ofrece en lugar del pueblo:

\begin{verso}{Éxodo 32:31–32}
Entonces volvió Moisés a \SENOR, y dijo: Te ruego, pues este pueblo ha cometido un gran pecado, porque se hicieron dioses de oro, que perdones ahora su pecado, y si no, ráeme ahora de tu libro que has escrito.
\end{verso}

La segunda es el pedido de ver la gloria (Éx 33:18). La tercera llega inmediatamente después de oír el Nombre:

\hebverso{אִם־נָא מָצָאתִי חֵן בְּעֵינֶיךָ אֲדֹנָי יֵלֶךְ־נָא אֲדֹנָי בְּקִרְבֵּנוּ \\ כִּי עַם־קְשֵׁה־עֹרֶף הוּא וְסָלַחְתָּ לַעֲוֹנֵנוּ וּלְחַטָּאתֵנוּ וּנְחַלְתָּנוּ}
\begin{verso}{Éxodo 34:9}
Y dijo: Si ahora, Señor, he hallado gracia en tus ojos, vaya ahora el Señor en medio de nosotros; porque es un pueblo de dura cerviz; y perdona nuestra iniquidad y nuestro pecado, y tómanos por tu heredad.
\end{verso}

\begin{filologia}{Éx 34:9}
\heb{חֵן} \emph{jen} & \heb{חנן} & ``gracia'': Moisés apela a la gracia que él mismo halló. \\
\heb{כִּי עַם־קְשֵׁה־עֹרֶף} \emph{ki am-qeshé óref} & \heb{קשׁה} & ``porque es un pueblo duro de cerviz''. La conjunción \emph{ki} puede leerse ``porque'' o ``aunque''; la RVR1960 dice ``porque''. \\
\heb{וְסָלַחְתָּ} \emph{vesalajtá} & \heb{סלח} & ``y perdonarás'': el verbo del perdón que solo tiene a Dios por sujeto en la Biblia hebrea. \\
\end{filologia}

Moisés usa como argumento lo que haría desistir a cualquiera: la dureza del pueblo. Si la misericordia dependiera de la calidad de Israel, no habría esperanza; precisamente porque el pueblo es duro, necesita un Dios ``tardo para la ira y grande en misericordia''. Dios responde: ``He aquí, yo hago pacto'' (Éx 34:10).

Moisés aprendió en el Sinaí a orar con las palabras que Dios había dicho de sí mismo. Cuando el pueblo se rebela otra vez tras el informe de los espías, cita el Nombre:

\hebverso{סְלַח־נָא לַעֲוֹן הָעָם הַזֶּה כְּגֹדֶל חַסְדֶּךָ}
\begin{verso}{Números 14:19–20}
Perdona ahora la iniquidad de este pueblo según la grandeza de tu misericordia, y como has perdonado a este pueblo desde Egipto hasta aquí. Entonces \SENOR{} dijo: Yo lo he perdonado conforme a tu dicho.
\end{verso}

Moisés no apela a la bondad de Israel, sino al carácter de Dios. Ese patrón (orar con el Nombre) es el que la tradición judía convirtió en el centro de su oración de perdón.

%===================================================================
\section{La misericordia en el judaísmo}

\subsection{Las trece medidas}
La tradición judía llama a Éx 34:6--7 las \emph{shelosh esré middot} (\heb{שְׁלֹשׁ עֶשְׂרֵה מִדּוֹת}), ``las trece medidas (o atributos) de misericordia''. La cuenta más conocida es la de Rabbenu Tam, que recogen los Tosafot al Talmud (\emph{Rosh Hashaná} 17b): los dos nombres iniciales son dos medidas distintas, la misericordia de Dios antes de que el hombre peque y después de que peque y se arrepienta.

\begin{center}\small\renewcommand{\arraystretch}{1.15}
\begin{tabular}{@{}cL{4.2cm}L{5.8cm}@{}}
\toprule
\textbf{N.º} & \textbf{Hebreo} & \textbf{Sentido} \\ \midrule
1 & \heb{יְהוָה} & \SENOR{} (misericordia antes del pecado) \\
2 & \heb{יְהוָה} & \SENOR{} (misericordia después del arrepentimiento) \\
3 & \heb{אֵל} & Dios, poderoso en misericordia \\
4 & \heb{רַחוּם} & compasivo \\
5 & \heb{וְחַנּוּן} & y clemente \\
6 & \heb{אֶרֶךְ אַפַּיִם} & lento para la ira \\
7 & \heb{וְרַב־חֶסֶד} & y abundante en \emph{jésed} \\
8 & \heb{וֶאֱמֶת} & y verdad \\
9 & \heb{נֹצֵר חֶסֶד לָאֲלָפִים} & que guarda el \emph{jésed} para millares \\
10 & \heb{נֹשֵׂא עָוֹן} & que carga la iniquidad (pecado intencional) \\
11 & \heb{וָפֶשַׁע} & y la rebelión \\
12 & \heb{וְחַטָּאָה} & y el pecado (sin intención) \\
13 & \heb{וְנַקֵּה} & y absuelve \\
\bottomrule
\end{tabular}\par
{\footnotesize\color{gris}\textit{Tabla 5.} Las trece medidas según Rabbenu Tam. \textit{Nota.} Tosafot a \emph{Rosh Hashaná} 17b (Sefaria).}
\end{center}

\precision{la cuenta no es única. Según los mismos Tosafot, Rabbenu Nissim no cuenta el primer nombre, porque hay una pausa entre los dos, y divide ``que guarda el \emph{jésed}'' y ``para millares''. La imagen del material de estudio sigue otra cuenta, que empieza en \emph{El}. La medida 13 (``y absuelve'') se obtiene cortando la frase \emph{venaqqé lo yenaqqé} (``absolver, no absuelve''): es una lectura rabínica, no el sentido llano del verso.}

\subsection{“Les mostró el orden de la oración”}
El Talmud comenta así Éx 34:6, ``Y pasando el \SENOR{} por delante de él, proclamó'':

\precision{lo que sigue es interpretación rabínica, tomada del Talmud de Babilonia (\emph{Rosh Hashaná} 17b) según la edición William Davidson en Sefaria (Steinsaltz, 2017).}

Rabí Yojanán dice que, si el verso no estuviera escrito, sería imposible decirlo: enseña que el Santo, bendito sea, se envolvió en un manto de oración como el que dirige la oración de la congregación, y le mostró a Moisés el orden de la oración. Le dijo:

\hebverso{כׇּל זְמַן שֶׁיִּשְׂרָאֵל חוֹטְאִין יַעֲשׂוּ לְפָנַי כַּסֵּדֶר הַזֶּה וַאֲנִי מוֹחֵל לָהֶם}

``Siempre que Israel peque, que hagan delante de mí conforme a este orden, y yo los perdonaré.'' Rav Yehudá añade:

\hebverso{בְּרִית כְּרוּתָה לִשְׁלֹשׁ עֶשְׂרֵה מִדּוֹת שֶׁאֵינָן חוֹזְרוֹת רֵיקָם}

``Hay un pacto hecho con las trece medidas: no vuelven vacías'', y lo prueba con Éx 34:10: ``He aquí, yo hago pacto''. Rashi, en su comentario a Éx 34:6, resume la primera idea: ``\SENOR, \SENOR'' es la medida de misericordia, una vez antes de que el hombre peque y otra después de que peque y se arrepienta.

La lógica es la misma de Moisés en Nm 14:18--19: el que pecó ora con las palabras que Dios dijo de sí mismo.

\subsection{La medida de misericordia y la medida del juicio}
Los sabios distinguen dos ``medidas'' en Dios: \heb{מִדַּת הָרַחֲמִים} (\emph{middat harajamim}), la medida de misericordia, unida al nombre \heb{יְהוָה}, y \heb{מִדַּת הַדִּין} (\emph{middat haddín}), la medida del juicio, unida a \emph{Elohim}. Rashi, comentando Gn 1:1, explica que primero ``surgió en el pensamiento'' de Dios crear el mundo con la medida del juicio; al ver que el mundo no podría subsistir, ``antepuso la medida de misericordia y la asoció a la medida del juicio''. En el Talmud, Dios es llamado con frecuencia \heb{רַחֲמָנָא} (\emph{Rajmaná}), ``el Misericordioso''.

El Targum Ónkelos, la traducción aramea de la Torá, vierte Éx 34:6--7 así:

\hebverso{וְאַעְבַּר יְיָ שְׁכִנְתֵּיהּ עַל אַפּוֹהִי וּקְרָא יְיָ יְיָ אֱלָהָא רַחֲמָנָא וְחַנָּנָא}

``Y el \SENOR{} hizo pasar su \emph{Shejiná} (presencia) delante de él, y proclamó: \SENOR, \SENOR, Dios misericordioso y clemente''. Y donde el hebreo dice ``y absolver, no absuelve'', Ónkelos traduce:

\hebverso{סָלַח לִדְתַיְבִין לְאוֹרַיְתֵיהּ וְלִדְלָא תַיְבִין לָא מְזַכֵּי}

``Perdona a los que vuelven a su Torá, y a los que no vuelven no los absuelve''.

\precision{son lecturas rabínicas (Targum Ónkelos según Sefaria). Ónkelos resuelve la tensión de Éx 34:7 con el arrepentimiento, igual que Pink: ``Dios muestra misericordia a los verdaderamente penitentes, pero no a los impenitentes'' (Pink, 2020, cap. 14).}

\subsection{En la oración y en el calendario}
Las trece medidas son el centro de las \emph{selijot} (\heb{סְלִיחוֹת}), las oraciones penitenciales que se recitan antes de Rosh Hashaná y en los días de ayuno (Elbogen, 1993). En Yom Kipur se repiten a lo largo del día, hasta la oración final (\emph{Neilá}). En la tarde de ese día se lee el libro de Jonás (\emph{Meguilá} 31a): la misericordia de Dios con Nínive, leída en el día del perdón de Israel. En Rosh Hashaná, en el rito del \emph{Tashlij}, se recita junto a aguas corrientes el final de Miqueas:

\hebverso{מִי־אֵל כָּמוֹךָ נֹשֵׂא עָוֹן וְעֹבֵר עַל־פֶּשַׁע לִשְׁאֵרִית נַחֲלָתוֹ \\ לֹא־הֶחֱזִיק לָעַד אַפּוֹ כִּי־חָפֵץ חֶסֶד הוּא}
\begin{verso}{Miqueas 7:18–19}
¿Qué Dios como tú, que perdona la maldad, y olvida el pecado del remanente de su heredad? No retuvo para siempre su enojo, porque se deleita en misericordia. El volverá a tener misericordia de nosotros; sepultará nuestras iniquidades, y echará en lo profundo del mar todos nuestros pecados.
\end{verso}

\begin{filologia}{Mi 7:18}
\heb{נֹשֵׂא עָוֹן} \emph{nosé avón} & \heb{נשׂא} & ``que carga la iniquidad'': cita literal de Éx 34:7. \\
\heb{כִּי־חָפֵץ חֶסֶד הוּא} \emph{ki jafets jésed hu} & \heb{חפץ} & ``porque él se deleita en \emph{jésed}'': la misericordia no es para Dios una concesión, sino su gusto. \\
\heb{מִי־אֵל כָּמוֹךָ} \emph{mi El kamoja} & --- & ``¿qué Dios como tú?''. Posible juego con el nombre del profeta, \emph{Mijá}, ``¿quién como (el \SENOR)?''. \\
\end{filologia}

En la oración diaria, la bendición que precede al \emph{Shemá} (\emph{Ahavá rabbá}) pide:
\hebverso{אָבִינוּ הָאָב הָרַחֲמָן הַמְרַחֵם רַחֵם עָלֵינוּ}
\noindent ``Padre nuestro, Padre misericordioso, que tienes compasión: ten compasión de nosotros''.

\subsection{“Sé semejante a él”}
El judaísmo saca de las medidas de Dios una ética. Abba Shaúl comenta ``este es mi Dios, y lo alabaré'' (Éx 15:2) con un juego de palabras sobre \heb{וְאַנְוֵהוּ} (\emph{ve'anvéhu}):

\hebverso{הֱוֵי דּוֹמֶה לוֹ מָה הוּא חַנּוּן וְרַחוּם אַף אַתָּה חַנּוּן וְרַחוּם}

``Sé semejante a él: como él es clemente y misericordioso, sé tú también clemente y misericordioso'' (\emph{Shabat} 133b). Rabí Jamá bar Janiná explica ``en pos del \SENOR{} vuestro Dios andaréis'' (Dt 13:4 [13:5]) como seguir las medidas del Santo: así como él vistió a los desnudos (Adán y Eva), visitó a los enfermos (Abraham), consoló a los que lloraban (Isaac) y sepultó a los muertos (Moisés), así también tú (\emph{Sotá} 14a).

\precision{interpretación rabínica (Talmud de Babilonia, edición William Davidson en Sefaria). Es el paralelo judío más cercano a ``Sed, pues, misericordiosos, como también vuestro Padre es misericordioso'' (Lc 6:36).}

\subsection{Qumrán}
El rollo de los Salmos de la cueva 11 de Qumrán (11QSal\textsuperscript{a}) incluye el Salmo 145, el acróstico que cita Éx 34:6 en su v.~8 (``Clemente y misericordioso es \SENOR, lento para la ira, y grande en misericordia''). En ese rollo, cada verso del salmo va seguido de un estribillo:

\hebverso{בָּרוּךְ יְהוָה וּבָרוּךְ שְׁמוֹ לְעוֹלָם וָעֶד}

``Bendito el \SENOR{} y bendito su nombre por siempre jamás''. La comunidad que copió el rollo cantaba el Nombre de Éx 34:6 bendiciendo el Nombre después de cada línea.

\precision{según la edición de Sanders (1965, DJD IV). Conviene cotejar el detalle con las imágenes del rollo en la Leon Levy Dead Sea Scrolls Digital Library antes de citarlo en público.}

%===================================================================
\section{La misericordia en Cristo, leída con Pablo}

\subsection{Pablo cita Éxodo 33:19}
Pablo, judío formado en las Escrituras de Israel, enfrenta en Romanos 9--11 la pregunta por la fidelidad de Dios a Israel. Y la responde con el texto del Sinaí:

\begin{verso}{Romanos 9:14–16}
¿Qué, pues, diremos? ¿Que hay injusticia en Dios? En ninguna manera. Pues a Moisés dice: Tendré misericordia del que yo tenga misericordia, y me compadeceré del que yo me compadezca. Así que no depende del que quiere, ni del que corre, sino de Dios que tiene misericordia.
\end{verso}

\begin{center}\small\renewcommand{\arraystretch}{1.25}
\begin{tabular}{@{}L{1.8cm}L{12.4cm}@{}}
\toprule
TM & \heb{וְחַנֹּתִי אֶת־אֲשֶׁר אָחֹן וְרִחַמְתִּי אֶת־אֲשֶׁר אֲרַחֵם} \\
LXX & \grc{ἐλεήσω ὃν ἂν ἐλεῶ καὶ οἰκτιρήσω ὃν ἂν οἰκτίρω} \\
Ro 9:15 & \grc{ἐλεήσω ὃν ἂν ἐλεῶ καὶ οἰκτιρήσω ὃν ἂν οἰκτίρω} \\
\bottomrule
\end{tabular}\par
{\footnotesize\color{gris}\textit{Tabla 6.} Éx 33:19 en el texto masorético, la LXX y Romanos 9:15.}
\end{center}

Pablo cita la LXX palabra por palabra. El contexto de Éx 33 es decisivo: Dios dice esa frase a Moisés justo después del becerro, cuando Israel no merecía nada. La libertad de la misericordia no es arbitrariedad; es la única razón por la que Israel siguió existiendo. ``No depende del que quiere, ni del que corre'' (\grc{οὐ τοῦ θέλοντος οὐδὲ τοῦ τρέχοντος}).

\subsection{Misericordia para todos}
El argumento de Romanos 9--11 empieza con Éx 33:19 y termina en doxología:

\begin{verso}{Romanos 11:32–33}
Porque Dios sujetó a todos en desobediencia, para tener misericordia de todos. ¡Oh profundidad de las riquezas de la sabiduría y de la ciencia de Dios! ¡Cuán insondables son sus juicios, e inescrutables sus caminos!
\end{verso}

\grc{ἵνα τοὺς πάντας ἐλεήσῃ}: ``para tener misericordia de todos'', judíos y gentiles. Los gentiles ``han alcanzado misericordia por la desobediencia'' de Israel, para que Israel también ``alcance misericordia'' (Ro 11:30--31). Más adelante Pablo resume la misión de Cristo con el par \emph{emet}--\emph{jésed}:

\begin{verso}{Romanos 15:8–9}
Pues os digo, que Cristo Jesús vino a ser siervo de la circuncisión para mostrar la verdad de Dios, para confirmar las promesas hechas a los padres, y para que los gentiles glorifiquen a Dios por su misericordia.
\end{verso}

Verdad para Jacob, misericordia para las naciones: es el mismo par de Mi 7:20 (``cumplirás la verdad a Jacob, y a Abraham la misericordia'').

\precision{la conexión entre Ro 15:8--9 y Mi 7:20 es una lectura propia; el par ``verdad''--``misericordia'' está en ambos textos.}

\subsection{“Padre de misericordias”}
\begin{verso}{2 Corintios 1:3}
Bendito sea el Dios y Padre de nuestro Señor Jesucristo, Padre de misericordias y Dios de toda consolación.
\end{verso}

\grc{ὁ πατὴρ τῶν οἰκτιρμῶν}: ``el Padre de las compasiones''. Pablo bendice a Dios con la forma de la bendición judía (\emph{barúj}, ``bendito'') y con un título cercano al de la oración de Israel: ``Padre misericordioso'' (\heb{הָאָב הָרַחֲמָן}). La palabra \emph{oiktirmós} es la que la LXX usa para \emph{rajamím}. En Efesios, el eco es el \emph{rav jésed} de Éx 34:6:

\begin{verso}{Efesios 2:4–5}
Pero Dios, que es rico en misericordia, por su gran amor con que nos amó, aun estando nosotros muertos en pecados, nos dio vida juntamente con Cristo (por gracia sois salvos).
\end{verso}

\grc{πλούσιος ὢν ἐν ἐλέει}, ``siendo rico en misericordia'', corresponde al \grc{πολυέλεος} de la LXX en Éx 34:6.

\precision{el paralelo entre 2~Co 1:3 y la oración judía es una lectura propia. La autoría paulina de Efesios es discutida; el título equivalente de 2~Co 1:3 está en una carta indiscutida.}

\subsection{Salvos por su misericordia}
\begin{verso}{Tito 3:4–5}
Pero cuando se manifestó la bondad de Dios nuestro Salvador, y su amor para con los hombres, nos salvó, no por obras de justicia que nosotros hubiéramos hecho, sino por su misericordia, por el lavamiento de la regeneración y por la renovación en el Espíritu Santo.
\end{verso}

\grc{κατὰ τὸ αὐτοῦ ἔλεος ἔσωσεν ἡμᾶς}: ``según su misericordia nos salvó''. Obras y misericordia se excluyen; como dice Pink, es una contradicción hablar de merecer misericordia. Pablo se presenta a sí mismo como ejemplo:

\begin{verso}{1 Timoteo 1:16}
Pero por esto fui recibido a misericordia, para que Jesucristo mostrase en mí el primero toda su clemencia, para ejemplo de los que habrían de creer en él para vida eterna.
\end{verso}

Pablo había sido ``blasfemo, perseguidor e injuriador'' (1~Ti 1:13). Como Moisés después del becerro, es un hombre que vio la gloria de Dios después de su pecado (Hch 9:3; 2~Co 4:6).

\precision{la autoría de las cartas pastorales es discutida; la misma enseñanza está en Ro 9:16 y Ef 2:4--5.}

\subsection{La gloria que Moisés pidió}
Moisés pidió ver la gloria y no pudo ver el rostro: ``No podrás ver mi rostro; porque no me verá hombre, y vivirá'' (Éx 33:20). Vio sus ``espaldas'' (33:23) y oyó su Nombre. Pablo compara su ministerio con el de Moisés (2~Co 3:7--4:6) y concluye:

\begin{verso}{2 Corintios 4:6}
Porque Dios, que mandó que de las tinieblas resplandeciese la luz, es el que resplandeció en nuestros corazones, para iluminación del conocimiento de la gloria de Dios en la faz de Jesucristo.
\end{verso}

En Cristo la gloria tiene rostro (\grc{πρόσωπον}), y ese rostro es misericordia. Juan lo dice con las palabras del Sinaí: ``vimos su gloria\ldots{} lleno de gracia y de verdad'' (Jn 1:14), y ``la ley por medio de Moisés fue dada, pero la gracia y la verdad vinieron por medio de Jesucristo'' (Jn 1:17). Es la tesis de esta obra: en Cristo ``habita corporalmente toda la plenitud de la Deidad'' (Col 2:9).

\precision{la conexión de 2~Co 4:6 con Jn 1:14--17 es de esta serie; Juan no es Pablo, pero ambos leen el Sinaí cumplido en Cristo.}

\subsection{El propiciatorio}
\begin{verso}{Romanos 3:25}
a quien Dios puso como propiciación por medio de la fe en su sangre, para manifestar su justicia, a causa de haber pasado por alto, en su paciencia, los pecados pasados.
\end{verso}

\grc{ἱλαστήριον} es en la LXX la tapa del arca, \heb{כַּפֹּרֶת} (\emph{kapporet}), el lugar donde Dios se encontraba con Israel: ``de allí me declararé a ti, y hablaré contigo de sobre el propiciatorio'' (Éx 25:22). Las tablas renovadas después de la revelación del Nombre se guardaron dentro del arca, bajo esa tapa (Dt 10:1--5). En la cruz, ``la misericordia y la verdad se encontraron'' (Sal 85:10): Dios perdona sin dejar de ser justo, ``a fin de que él sea el justo, y el que justifica'' (Ro 3:26).

\precision{la imagen de la ley guardada bajo la tapa de la misericordia es una lectura tipológica propia.}

%===================================================================
\section{Aplicación: un atributo comunicable}

\subsection{“Sed misericordiosos”}
\begin{verso}{Lucas 6:36}
Sed, pues, misericordiosos, como también vuestro Padre es misericordioso.
\end{verso}
\begin{verso}{Mateo 5:7}
Bienaventurados los misericordiosos, porque ellos alcanzarán misericordia.
\end{verso}

\grc{γίνεσθε οἰκτίρμονες καθὼς ὁ πατὴρ ὑμῶν οἰκτίρμων ἐστίν}: \grc{οἰκτίρμων} es la palabra con que la LXX traduce \emph{rajum} en Éx 34:6. Jesús pide a sus discípulos el primer rasgo del Nombre. La misericordia es comunicable: el ser humano puede reflejarla, no producirla. Se recibe primero y se da después; ``alcanzarán misericordia'' (\grc{ἐλεηθήσονται}) es un pasivo: Dios la da. Miqueas lo resumió para Israel:

\hebverso{עֲשׂוֹת מִשְׁפָּט וְאַהֲבַת חֶסֶד וְהַצְנֵעַ לֶכֶת עִם־אֱלֹהֶיךָ}
\begin{verso}{Miqueas 6:8}
Oh hombre, él te ha declarado lo que es bueno, y qué pide \SENOR{} de ti: solamente hacer justicia, y amar misericordia, y humillarte ante tu Dios.
\end{verso}

\subsection{¿Quién no alcanzará misericordia?}
\begin{verso}{Santiago 2:13}
Porque juicio sin misericordia se hará con aquel que no hiciere misericordia; y la misericordia triunfa sobre el juicio.
\end{verso}

\begin{filologia}{Stg 2:13}
\grc{ἀνέλεος} \emph{anéleos} & \grc{ἔλεος} & ``sin misericordia'': el juicio sin \emph{éleos}. \\
\grc{κατακαυχᾶται} \emph{katakaujátai} & \grc{καυχάομαι} & ``se gloría por encima de, triunfa sobre'': la misericordia tiene la última palabra sobre el juicio. \\
\end{filologia}

Aquí se afirma claramente que el que no muestra misericordia será juzgado sin ella. Es lo que enseña la parábola del siervo: ``¿No debías tú también tener misericordia de tu consiervo, como yo tuve misericordia de ti?'' (Mt 18:33). Diez mil talentos perdonados, cien denarios exigidos (Mt 18:24, 28). Santiago, judío, dice lo mismo que Éx 34:7: la misericordia es abundante, pero nadie se burla de ella.

\precision{santiago no es de Pablo; el texto se incluye porque forma parte del material de estudio de la serie.}

\subsection{“Ve, y haz tú lo mismo”}
\begin{verso}{Lucas 10:33–37}
Pero un samaritano, que iba de camino, vino cerca de él, y viéndole, fue movido a misericordia; y acercándose, vendó sus heridas, echándoles aceite y vino; y poniéndole en su cabalgadura, lo llevó al mesón, y cuidó de él. [\ldots] ¿Quién, pues, de estos tres te parece que fue el prójimo del que cayó en manos de los ladrones? El dijo: El que usó de misericordia con él. Entonces Jesús le dijo: Ve, y haz tú lo mismo.
\end{verso}

\begin{filologia}{Lc 10:33, 37}
\grc{ἐσπλαγχνίσθη} \emph{esplanjnísthē} & \grc{σπλάγχνα} & ``se conmovió en las entrañas'' (10:33): \emph{rajamím}. \\
\grc{ὁ ποιήσας τὸ ἔλεος} & \grc{ποιέω} & ``el que hizo misericordia'' (10:37): la expresión hebrea \heb{עָשָׂה חֶסֶד} (\emph{asá jésed}), ``hacer \emph{jésed}''. \\
\end{filologia}

Un intérprete de la ley pregunta ``¿y quién es mi prójimo?'' (Lc 10:29). Jesús da vuelta la pregunta: no es a quién debo amar, sino de quién me hago prójimo. La misericordia se siente (entrañas) y se hace (\emph{asá jésed}).

\subsection{Siquem: el pacto sin \emph{jésed}}\label{sec:siquem}
El material de estudio termina con una advertencia desde Siquem, el lugar donde Abram levantó su primer altar:

\hebverso{וַיָּשִׂימוּ לָהֶם בַּעַל בְּרִית לֵאלֹהִים \\ וְלֹא־עָשׂוּ חֶסֶד עִם־בֵּית יְרֻבַּעַל גִּדְעוֹן}
\begin{verso}{Jueces 8:33–35}
Pero aconteció que cuando murió Gedeón, los hijos de Israel volvieron a prostituirse yendo tras los baales, y escogieron por dios a Baal-berit. Y no se acordaron los hijos de Israel de \SENOR{} su Dios, que los había librado de todos sus enemigos en derredor; ni se mostraron agradecidos con la casa de Jerobaal, el cual es Gedeón, conforme a todo el bien que él había hecho a Israel.
\end{verso}

\begin{filologia}{Jue 8:33--35; 9:4, 46}
\heb{בַּעַל בְּרִית} \emph{Ba'al berit} & --- & ``señor del pacto'': un dios cananeo del pacto, con templo en Siquem (Jue 9:4). \\
\heb{אֵל בְּרִית} \emph{El berit} & --- & ``dios del pacto'' (Jue 9:46): el templo-fortaleza de la torre de Siquem. \\
\heb{וְלֹא־עָשׂוּ חֶסֶד} \emph{velo-asú jésed} & \heb{חסד} & ``y no hicieron \emph{jésed}'' (RVR1960: ``ni se mostraron agradecidos''). \\
\end{filologia}

Donde Abram había dicho ``no'' a los dioses de la tierra (Gn 12:6--7), Israel cambió al Dios del pacto por un ``señor del pacto''. El texto une dos olvidos: no se acordaron del \SENOR{} (8:34) y no hicieron \emph{jésed} con la casa de Gedeón (8:35). Quien olvida el \emph{jésed} de Dios deja de hacer \emph{jésed} a los demás. Es otra vez la gloria cambiada.

\precision{el vínculo entre 8:34 y 8:35 está en el texto. Leer Baal-berit como una imitación cananea del pacto de Israel es una interpretación (Block, 1999).}

\subsection{Cuatro respuestas}
\begin{enumerate}
\item \textbf{Elegir su gloria.} No cambiar la gloria del Dios incorruptible por ídolos de hoy (Ro 1:23), sino decir el ``no'' de Abram en Siquem (Gn 12:6--7).
\item \textbf{Orar con su Nombre.} Después de pecar, apelar al carácter de Dios y no a nuestros méritos, como Moisés (Éx 34:9; Nm 14:19) y como la oración de Israel.
\item \textbf{Recibir y dar.} Perdonar como fuimos perdonados (Mt 18:33) y ser misericordiosos como el Padre (Lc 6:36).
\item \textbf{Hacerse prójimo.} Sentir con las entrañas y actuar (Lc 10:37; Mi 6:8).
\end{enumerate}

\begin{verso}{Romanos 12:1}
Así que, hermanos, os ruego por las misericordias de Dios, que presentéis vuestros cuerpos en sacrificio vivo, santo, agradable a Dios, que es vuestro culto racional.
\end{verso}

\grc{διὰ τῶν οἰκτιρμῶν τοῦ θεοῦ}: toda la ética de Romanos 12--15 se apoya en las misericordias expuestas en Romanos 1--11.

%===================================================================
\section{Síntesis}

\begin{tcolorbox}[colback=papiro,colframe=azul,boxrule=1pt,arc=3pt]
Cuando Israel cambió su gloria por un becerro, Dios mostró la suya: ``¡\SENOR! ¡\SENOR! fuerte, misericordioso y piadoso'' (Éx 34:6). Su misericordia es gracia libre (\emph{janán}), amor leal (\emph{jésed}), ternura de entrañas (\emph{rajamím}) y clemencia que saca del juicio (\emph{jamal}), sin dejar de ser justa. Israel la convirtió en su oración; Pablo la vio en el rostro de Cristo (2~Co 4:6), ``Padre de misericordias'' (2~Co 1:3). Quien la recibe, la refleja (Lc 6:36).
\end{tcolorbox}

\begin{center}\small\renewcommand{\arraystretch}{1.3}
\begin{tabular}{@{}L{3.0cm}L{11.4cm}@{}}
\toprule
\textbf{Mirada} & \textbf{Aporte al atributo} \\ \midrule
Hebrea & Cuatro palabras (\emph{janán}, \emph{jésed}, \emph{rajamím}, \emph{jamal}); el Nombre de Éx 34:6--7 y sus ecos; Moisés intercede con el Nombre; Lot, Efraín, Nínive; Baal-berit. \\
Judía & Las trece medidas; ``les mostró el orden de la oración''; \emph{middat harajamim}; \emph{Rajmaná}; Ónkelos; \emph{selijot} y Yom Kipur; ``sé semejante a él''. \\
Pink & Nace de la bondad y presupone la miseria; general, especial y soberana; su ejercicio lo regula la voluntad de Dios; no se presume. \\
Paulina & Ro 9:15 = Éx 33:19; misericordia para todos (Ro 11:32); Padre de misericordias; salvos por su misericordia; la gloria en el rostro de Cristo. \\
\bottomrule
\end{tabular}\end{center}

\begin{verso}{Jeremías 9:24 [9:23]}
Mas alábese en esto el que se hubiere de alabar: en entenderme y conocerme, que yo soy \SENOR, que hago misericordia, juicio y justicia en la tierra; porque estas cosas quiero, dice \SENOR.
\end{verso}

El verso que une esta obra pone la misericordia primero: \heb{עֹשֶׂה חֶסֶד מִשְׁפָּט וּצְדָקָה}, ``que hago \emph{jésed}, juicio y justicia''. Conocer a Dios es conocer al que hace \emph{jésed}.

%===================================================================
\section*{Glosario del capítulo}
\begin{description}\setlength{\itemsep}{2pt}
\item[\emph{Anéleos}] (\grc{ἀνέλεος}). Sin misericordia (Stg 2:13).
\item[\emph{Éleos}] (\grc{ἔλεος}). Misericordia, compasión activa; en la LXX traduce \emph{jésed}.
\item[\emph{Eleēmōn}] (\grc{ἐλεήμων}). Misericordioso (Mt 5:7; He 2:17).
\item[\emph{Emet}] (\heb{אֱמֶת}). Verdad, firmeza, fidelidad.
\item[\emph{Eúsplanjnos}] (\grc{εὔσπλαγχνος}). ``De buenas entrañas'', misericordioso (Ef 4:32).
\item[\emph{Hósios}] (\grc{ὅσιος}). Leal, santo; \emph{ta hósia Dauíd}, las lealtades prometidas a David (Hch 13:34).
\item[\emph{Janán}] (\heb{חָנַן}). Inclinarse con favor hacia un inferior; tener gracia.
\item[\emph{Jamal}] (\heb{חָמַל}). Compadecerse; perdonar la vida.
\item[\emph{Jasid}] (\heb{חָסִיד}). Misericordioso: el hombre de \emph{jésed} (Mi 7:2; Jer 3:12).
\item[\emph{Jésed}] (\heb{חֶסֶד}). Amor leal del pacto.
\item[\emph{Jus}] (\heb{חוּס}). Tener lástima, perdonar la vida (Jon 4:11).
\item[\emph{Kavod}] (\heb{כָּבוֹד}). Gloria; literalmente, ``peso''.
\item[\emph{Middat harajamim}] (\heb{מִדַּת הָרַחֲמִים}). La medida de misericordia, en la teología rabínica.
\item[\emph{Nijam}] (\heb{נִחַם}). Conmoverse, cambiar el curso (Jue 2:18).
\item[\emph{Oiktirmós}] (\grc{οἰκτιρμός}). Compasión (2~Co 1:3; Ro 12:1).
\item[\emph{Polýsplanjnos}] (\grc{πολύσπλαγχνος}). ``De muchas entrañas'' (Stg 5:11).
\item[\emph{Rajamím}] (\heb{רַחֲמִים}). Entrañas; compasión.
\item[\emph{Rajmaná}] (\heb{רַחֲמָנָא}). ``El Misericordioso'', nombre de Dios en el Talmud.
\item[\emph{Rajum}] (\heb{רַחוּם}). Compasivo (Éx 34:6).
\item[\emph{Réjem}] (\heb{רֶחֶם}). Vientre, matriz.
\item[\emph{Selijot}] (\heb{סְלִיחוֹת}). Oraciones penitenciales judías centradas en las trece medidas.
\item[\emph{Shelosh esré middot}] (\heb{שְׁלֹשׁ עֶשְׂרֵה מִדּוֹת}). Las trece medidas de misericordia de Éx 34:6--7.
\item[\emph{Splánjna}] (\grc{σπλάγχνα}). Entrañas; ``entrañable misericordia'' (Lc 1:78; Col 3:12).
\item[\emph{Splanjnízomai}] (\grc{σπλαγχνίζομαι}). Conmoverse en las entrañas.
\item[\emph{Tejinná}, \emph{tajanunim}] (\heb{תְּחִנָּה}, \heb{תַּחֲנוּנִים}). Súplica; la gracia que se pide (Esd 9:8; Jer 31:9).
\item[\emph{Tov ayin}] (\heb{טוֹב־עַיִן}). ``Bueno de ojo'': generoso, misericordioso (Pr 22:9).
\end{description}

